\section{Related Work}

\paragraph{RL for LLM reasoning and verifiable rewards.}
Reinforcement learning has become a central post-training tool for improving mathematical reasoning, code generation, and other tasks where outcomes can be checked automatically. Systems such as DAPO scale reinforcement learning for \llm{} reasoning with mechanisms aimed at stable optimization and entropy control \cite{dapo_2025}. Recent \rlvr{} work studies why verifier-based training can improve reasoning, how spurious or clipped rewards interact with entropy, and how exploration changes the effectiveness of group-relative updates \cite{exploration_exploitation_rlvr_2026,wen2025reinforcement,he2025random}. Our work is complementary: rather than proposing a new optimizer, we formulate verifier feedback as a low-bandwidth information channel and measure its usable signal through reward variance and reward--score covariance.

\paragraph{Entropy, exploration, and zero-variance prompts.}
Several recent methods identify entropy collapse, insufficient rollout diversity, or zero-variance prompt groups as bottlenecks in \llm{} \rlvr{}. DAPO emphasizes entropy-stabilized large-scale RL, while follow-up work studies selective entropy regularization, bidirectional entropy modulation, entropy-curve control, and quantile-based advantage estimation \cite{dapo_2025,siren2025,asymgrpo2026,entrocraft2026,qae2025}. Other methods increase trajectory-level diversity or explicitly handle zero-variance prompts \cite{lookahead_rollouts_2026,eepo_2026,no_prompt_left_behind_2026}. These methods can be viewed through our lens as attempts to increase exposed verifier variance, keep prompts near the decision boundary, or improve the conversion from reward variation to policy updates.

\paragraph{Sparse, dense, and hybrid rewards.}
Verifiable rewards are attractive because they are low-bias with respect to precise task specifications, but they often compress rich trajectories into terminal binary outcomes. Hybrid and dense-reward methods attempt to recover part of the missing feedback bandwidth by combining sparse verifiers with learned reward models or graded supervisory signals \cite{hybrid_reinforcement_2026}. Our formulation separates this issue into reward-model bias, evaluator noise, reward variance, and update alignment. This distinction is important because denser feedback is useful only when it adds policy-controllable, update-aligned variance rather than merely increasing noisy reward dispersion.

\paragraph{Self-play and renewable feedback.}
AlphaZero-like self-play demonstrates a different way to overcome sparse terminal rewards: the reward itself may remain binary and terminal, but the rollout distribution is renewed by playing against an opponent of comparable strength \cite{silver2017mastering}. This keeps outcomes near a high-variance frontier and makes the reward-information integral grow with continued self-play. We use this regime as a contrast to fixed-dataset \llm{} \rlvr{}, where the prompt-verifier source is fixed and many prompts eventually become all-correct or all-wrong.

\paragraph{Diffusion and generative-model RL.}
Diffusion models introduce long generation trajectories and delayed or multi-objective feedback \cite{ho2020denoising}. In this setting, the main bottleneck is not only whether final-sample rewards have variance, but also whether those rewards have useful covariance with intermediate denoising decisions. The diffusion section of this paper is currently left as a placeholder, but the common covariance formulation is designed to extend to image and video generation rewards, including aesthetics, prompt alignment, temporal consistency, and safety.
